% Source: source_md/intelligence_part2_combined_ja.md
\section{応答チャネルと沈黙世界 — 閾値現象としての C-幾何学}


\S{}3 で示した通り、推論層 $L_F$ は内部停止条件を持たない。本章では、この非閉包性が推論の操作環境によってどのように現れるかを分析する。推論環境は応答的世界と沈黙的世界に大別されるが、\S{}2.6 で予告した通り、この区別は本体論的範疇ではなく閾値現象である。本章末尾では、沈黙世界における無限後退の運動がログとして堆積し、ログ層 $L_R$ の予兆を生むことを示す。これが \S{}5(第二中心)への助走となる。

\bigskip


\subsection{応答的世界と沈黙的世界}


推論主体が環境に対して推論出力(問い・仮説・行為)を投げかけたとき、環境がそれに応答を返すか否かによって、二つの様態が区別される。

\textbf{応答的世界}:
推論出力に対して環境が応答を返す世界。C-幾何学の語彙では、$\nabla C$ が双方向に閉ループを形成する領域である。推論主体の出力が環境の状態を変え、その変化が主体に戻り、次の推論を駆動する。他者との対話、実験、社会的相互作用がこれにあたる。

\textbf{沈黙的世界}:
推論出力に対して環境が応答を返さない世界。$\nabla C$ が一方向のみで、ループが閉じない領域である。推論主体は出力できるが、それに対する応答が返らない。自然現象の究極的理由、過去の確定した事象、未来の偶然性などがこれにあたる。

形式的に述べれば、応答的世界では推論主体ノード $i$ と環境ノード $j$ の間に双方向の相互拘束 $H_{ij}, H_{ji}$ がともに有意な値を持ち、閉包度 $C_{ij}$ が動的に更新される。沈黙的世界では $H_{ji} \approx 0$(環境から主体への拘束が返らない)であり、$C_{ij}$ は主体側からの一方的な更新に留まる。

\bigskip


\subsection{沈黙性は欠陥ではない}


ここで強調すべきは、沈黙性が環境の欠陥や不完全性ではないという点である。

自然現象は、正しく問いかけても、その究極的理由については沈黙する。「なぜ物理定数はこの値なのか」という問いに、自然は応答を返さない。過去の確定した事象は、それがなぜそうであったかについて、もはや応答しない。未来の偶然性は、それが確定する前には応答を持たない。

これらの沈黙は、知識の不足によるものではなく、構造的なものである。応答チャネルがそもそも存在しない、あるいは閉じている領域が存在する。沈黙的世界は、応答的世界の「まだ解明されていない部分」ではなく、応答ループが原理的に閉じない領域として区別される。

この点は重要である。なぜなら、沈黙的世界に対して応答的世界と同じ推論戦略を適用すると、後述する無限後退に陥るからである。

\bigskip


\subsection{応答的/沈黙的は閾値現象である}


\S{}2.6 の主張 2.2 を、本章の文脈で再確認する。

応答的世界と沈黙的世界は、二種類の世界が別個に存在することを意味しない。同一の環境が、相互拘束 $H_{ij}$ のパラメータ位置と、推論主体の $(v_{\text{inf}}, r_{\text{res}}, \eta)$ の組によって、応答的にも沈黙的にも現れる。

具体例で述べる。ある自然現象は、適切な時間スケールと解像度で問いかければ応答を返す(実験的に検証可能な因果関係として)。しかし、過度に高速・高解像度で問いかけると、応答が間に合わず、あるいは応答が解像度の細部に追いつかず、同じ現象が沈黙的に見える。逆に、適切な閾値設定によって、かつて沈黙的だった領域が応答的になることもある。これが科学の歴史において繰り返されてきたことである。

したがって、応答的/沈黙的の境界は固定されておらず、相互拘束パラメータ上を連続的に移動する閾値である。そして \S{}2.7 で導入した閾値移動能力 $\mu_\theta$ は、この境界を状況に応じて動かす能力にほかならない。

この観察は本章の後半で決定的になる。無限後退は「沈黙的世界に固定された主体」に起こるのであって、閾値を動かせる主体は無限後退を回避する経路を持つ。

\begin{figure}[H]
\centering
\includegraphics[width=\linewidth]{fig2_responsive_silent_threshold_ja.pdf}
\caption{応答的世界と沈黙的世界。両者は世界の存在論的な種別ではなく、$\mu_\theta$ が動かす閾値境界の左右にどちらが落ちるかで決まる相である。}
\label{fig:part2-responsive-silent}
\figurenote{左右は同一の環境・同一の $C$-地景を共有する。橙の閾値境界が動くことで、同じ世界が応答的にも沈黙的にも現れる。}{応答的相では外部相互拘束が局所停止を供給する。沈黙的相では停止供給が断たれ、停止を別の場所へ外部化する圧力が生じる。}
\end{figure}

\bigskip


\subsection{沈黙世界における無限後退}


沈黙的世界に固定された推論主体——すなわち閾値を動かせず、環境を沈黙的として扱い続ける主体——を考える。このとき、以下の三条件が同時に成立する。

\textbf{条件 (1)}: 環境が応答を返さない(沈黙性、$H_{ji} \approx 0$)。

\textbf{条件 (2)}: 正当化は推論によってのみ生成される(\S{}3 の前提 P2)。

\textbf{条件 (3)}: 停止は推論から導出できない(\S{}3 の定理 3.1、系 3.2)。

これら三条件が同時に成立すると、推論は無限後退に陥る。

論証は次の通りである。条件 (1) により、推論を停止させる外部応答は来ない。条件 (2) により、正当化を進めるには推論を続けるしかない。条件 (3) により、推論の内部から停止条件を導くこともできない。したがって推論は停止する手段を持たず、無限に後退する。

これは \S{}3.5 で論じた正当化の無限後退の、操作環境における具体的現れである。\S{}3 では抽象的に「内部停止述語は無限後退を生む」と述べたが、\S{}4 ではそれが「沈黙的世界に固定された主体において現実に発生する」ことが示される。

\bigskip


\subsection{無限後退と分岐爆発の相同}


\S{}4.4 の無限後退は、Part I \S{}7 で導入された分岐爆発と構造的に相同である。

無限後退は、推論軌道が停止せず延々と続く現象として記述される。分岐爆発は、推論が複数の方向に分岐し、各分岐がさらに分岐して指数的に増殖する現象として記述される。両者は \S{}3.7 で示した通り、内部不動点不在の二つの側面である:

\begin{itemize}
\item 内部不動点不在 → 軌道は静止しない(無限後退)
\item 内部不動点不在 + 確率構造(性質 F2) → 軌道は単一でない(分岐爆発)
\end{itemize}


沈黙的世界では、この二つが結合する。推論は停止せず(無限後退)、かつ単一の経路に収束せず(分岐爆発)、二重に不安定な動力学を示す。これが、沈黙的世界に固定された主体が経験する推論の発散である。

この発散は、\S{}3.8 で述べた三つの病的引力子(過剰確信・麻痺・妄想的閉包)へと主体を導く。停止できない推論は、早すぎる停止(過剰確信)、無限懐疑(麻痺)、あるいは自己封印的合理化(妄想的閉包)のいずれかに落ちる。これらはすべて $L_E$ を欠いた $L_F$ の引力子である。

\bigskip


\subsection{ログ層 $L_R$ の予兆}


ここで、本論文の第二中心 \S{}5 への決定的な橋渡しが現れる。

沈黙的世界における推論は停止しない。しかし、停止しない推論の運動は、痕跡を残す。推論軌道は状態空間を通過し、その通過の記録——どの状態を経たか、どの分岐を試みたか、どこで発散したか——が堆積する。

この堆積した痕跡を \textbf{ログ} と呼ぶ。そして、ログに対する再帰的な参照運動——過去の推論軌跡を読み返し、パターンを抽出し、現在の推論に反映する運動——が生じる層を、\textbf{ログ層} $L_R$ と呼ぶ(正式な定義は \S{}5.2)。

重要なのは、$L_R$ が $L_F$ の運動から自然に生じるという点である。$L_F$ が停止せず運動し続ける限り、その痕跡は堆積し、再帰参照の対象となる。すなわち、$L_F$ の非閉包性そのものが、$L_R$ を生み出す。停止できないがゆえに痕跡が溜まり、痕跡が溜まるがゆえに再帰参照が可能になる。

これは逆説的な構造である。$L_F$ の限界(停止できないこと)が、$L_F$ を超える層($L_R$)の生成条件になっている。そして \S{}5 で示すように、$L_R$ の再帰運動と外部相互拘束 $H_{ij}$ の関数として、創発停止層 $L_E$ が生じる。$L_E$ こそが、$L_F$ が単独では持ちえなかった停止を、個体内に外部依存的な形でもたらす。

すなわち、無限後退は単なる行き詰まりではない。それは $L_R$ を生成し、$L_R$ を通じて $L_E$ の創発を準備する、動力学的に生産的な運動でもある。

\bigskip


\subsection{次章への橋渡し}


本章で確立したのは次の経路である:

\[
L_F \text{の非閉包(§3)} \;\longrightarrow\; \text{沈黙世界での無限後退(§4.4)} \;\longrightarrow\; \text{ログの堆積} \;\longrightarrow\; L_R \text{の予兆(§4.6)}
\]

\S{}5 では、この $L_R$ から創発停止層 $L_E$ がいかにして生じるかを展開する。具体的には、ログ層 $L_R$ の定義、$L_E$ の創発機構(閾値移動 $\mu_\theta$ による (B) 組み替えと、相互拘束 $H_{ij}$ に沿った停止演算子設置領域の投影)、外部停止演算子 $E$ の構造、そして外部参照場 $\Phi_{\text{ext}}$ が個体間 $L_E$ ネットワーク $\Phi_{\text{net}}$ と同型であることが示される。

\S{}3 が「$L_F$ は閉じない」を、\S{}5 が「それでも $L_E$ を通じて停止する」を示す。本章 \S{}4 は両者を繋ぐ蝶番であり、特に無限後退が $L_R$ を生むという発見が、第一中心と第二中心を動力学的に連結する。

\bigskip


\subsection{要約}


本章で示したのは次の点である:

\begin{enumerate}
\item 推論環境は応答的世界($\nabla C$ 双方向閉ループ)と沈黙的世界($\nabla C$ 一方向)に区別される
\item 沈黙性は欠陥ではなく構造的なものである
\item しかし応答的/沈黙的は本体論的範疇ではなく、相互拘束パラメータ上の閾値現象である(\S{}2.6)
\item 沈黙世界に固定された主体では、三条件(沈黙・推論依存・停止不能)が無限後退を生む
\item 無限後退は分岐爆発と相同であり、三つの病的引力子へ導く
\item しかし無限後退の運動はログを堆積させ、ログ層 $L_R$ の予兆を生む
\item $L_F$ の非閉包性そのものが $L_R$ の生成条件であり、これが \S{}5 の $L_E$ 創発を準備する
\end{enumerate}


それでは \S{}5 において、本論文の第二中心——創発停止層 $L_E$ の動力学——を展開する。

\bigskip
